% Generated draft: requires mathematical and source review.
\documentclass[11pt]{article}
\usepackage[margin=25mm]{geometry}
\usepackage{fontspec}
\setmainfont{Latin Modern Roman}
\usepackage{amsmath,amssymb,mathtools,unicode-math}
\setmathfont{Latin Modern Math}
\usepackage{enumitem,needspace,xurl,hyperref}
\setlength{\parindent}{0pt}
\setlength{\parskip}{5pt}
\setlength{\emergencystretch}{3em}
\newcommand{\N}{\mathbb{N}}
\newcommand{\Z}{\mathbb{Z}}
\newcommand{\Q}{\mathbb{Q}}
\newcommand{\R}{\mathbb{R}}
\newcommand{\C}{\mathbb{C}}
\newcommand{\F}{\mathbb{F}}
\newcommand{\A}{\mathbb{A}}
\newcommand{\PP}{\mathbb{P}}
\newcommand{\E}{\mathbb{E}}
\newcommand{\eps}{\varepsilon}
\newcommand{\dd}{\,\mathrm{d}}
\begin{document}
\section*{UnsolvedMath ID: 6700063}
\textbf{Source problem code:} catalogue entry
\subsection{Definitions and mathematical statement}
For a smooth Riemannian metric $g$, the classical condition $\operatorname{Sc}_g\geq0$ is a pointwise inequality involving second derivatives of the metric. A continuous Riemannian metric is a continuous positive-definite symmetric tensor; a merely bounded measurable tensor has still less regularity, and neither class has classical pointwise scalar curvature defined by the smooth formula.

Develop an appropriate extension of the condition $\operatorname{Sc}(X)\geq0$ to some singular spaces, in particular manifolds with continuous metrics. The supplied sentence tentatively continues ``bounded measurable?'' and then stops, so bounded-measurable metrics are a possible proposed extension, not an explicitly completed universal domain.

A candidate theory should specify its admissible objects and the meaning of nonnegative scalar curvature, including how it relates to the classical condition on smooth metrics. Which stability, locality, or global comparison/rigidity principles are required remains unspecified. For measurable metrics the representative convention and any ellipticity bounds must also be fixed; they are not added here. The record therefore poses a definition-building program, not a supplied synthetic inequality or a demand that one condition work on every singular metric space.
\subsection{Short English statement}
Extend nonnegative scalar curvature to suitable singular spaces or continuous metrics, with bounded measurable metrics mentioned only tentatively. Specify the object class and the properties the new definition should preserve.
\subsection{Sources}
\begin{itemize}
\item Misha Gromov, \emph{101 Questions, Problems and Conjectures around Scalar Curvature} (2017), printed tag [?67], occurrence 66, PDF page 61, \url{https://www.ihes.fr/~gromov/wp-content/uploads/2018/08/101-problemsOct1-2017.pdf}; accessed 2026-07-29, PDF-text-normalized extraction. The supplied sentence is truncated and the provenance describes the public author PDF as incomplete.
\item The research summary mentions volumically positive scalar curvature, $C^0$-metric scalar-curvature frameworks, and later work of Gromov and others, without individual publication identifiers or exact axioms.
\item Local registry \texttt{data/unsolvedmath/problems.json}, permanent UnsolvedMath ID $6700063$, source code catalogue entry, \emph{Scalar Curvature Question [?67]}. Dataset attribution: Ulam AI and the UnsolvedMath Contributors.
\end{itemize}
\end{document}
